\begin{table}[htbp]\centering
\caption{Complete tolerance-axis partition of the original scalar frontiers}\label{tab:frontier47}
\begin{tabular}{lr}\toprule Relation on a partition interval & Number\\\midrule
Both attain; witness uses fewer queries & 19 \\
Both attain; spline uses fewer queries & 1 \\
Both attain; same prefix queries & 20 \\
Only witness attains within the cap & 4 \\
Only spline attains within the cap & 0 \\
Both exhaust the cap & 4 \\
\bottomrule\end{tabular}\end{table}

The complete curves contain 19 nonempty tolerance intervals on which both methods attain and the witness uses fewer Bellman queries, one favoring the spline, and twenty with identical counts. Four further intervals are attainable only by the witness; four below both final bounds are unattained. Interval counts are not probabilities or measures of economic importance. Above the largest breakpoint in each cell both first-rung counts coincide. The original targets remain unchanged and still have identical crossing resolutions. These descriptive intervals identify where a sharper certificate changes a resource decision without manufacturing new prospective targets.
